
\documentclass[11pt]{article}
\usepackage[margin=1in]{geometry}
\usepackage[T1]{fontenc}
\usepackage[utf8]{inputenc}
\usepackage{lmodern}
\usepackage{amsmath}
\usepackage{amssymb}
\usepackage{graphicx}
\usepackage{booktabs}
\usepackage{float}
\usepackage{hyperref}
\hypersetup{hidelinks}
\setlength{\parskip}{0.5em}
\setlength{\parindent}{0pt}
\begin{document}
\sloppy
\section{Step 4: Primary ROI-Based Statistical Analysis of Abstract versus Concrete Questions}
\label{sec:step4_fnirs_analysis}
% Analyst note added 2026-03-20: the descriptive waveform plots in the implemented Step 4 pipeline use a fixed onset-locked window from -2 s to +20 s and participant-weighted grand means on a 0.1 s common grid. This choice affects only the descriptive waveform figures and does not change the primary inferential ROI statistics, which still use the exact question_start-to-answer trial window specified below.
% Analyst note added 2026-03-20: by explicit user request, PID025 / 2025-11-25_002 and PID034 / 2025-11-27_003 are manually included in Step 4 despite Step 3 REVIEW_BEFORE_STEP4 status. This is a deliberate deviation from the written participant inclusion rule below and is recorded in the Step 4 inclusion table and report.
% Analyst note added 2026-03-21: the separate Step 4 visualization report uses a participant digitization point cloud plus source/detector/channel geometry in place of an fsaverage brain-surface rendering, because the local environment does not provide the required pyvista/fsaverage stack for direct reproduction of the example code. The onset-locked topographic figures in that supplementary report use a fixed epoch window from -3.5 s to +13.2 s with a -2 s to 0 s baseline for descriptive visualization only.

\subsection{Objective}
The objective of Step~4 is to perform a simple, pre-specified, and statistically interpretable
analysis of the fNIRS signal to test whether abstract and concrete questions differ in the
left-lateralized montage used in this study.

This step should remain deliberately simple.
It should provide:
\begin{itemize}
    \item one primary inferential analysis;
    \item exact p-values;
    \item confidence intervals;
    \item effect sizes;
    \item a short conclusion written in clear language.
\end{itemize}

More complex modelling, such as mixed-effects models, full event-duration GLMs, correctness interactions,
or trial-level hierarchical analyses, can be deferred to Step~5.

\subsection{Primary Scientific Question}
The primary scientific question is:
\begin{quote}
Is the difference between abstract and concrete questions different in the \textbf{front} versus the
\textbf{back} of the left lateral montage?
\end{quote}

This question is more appropriate than averaging all channels together, because the expected effect is
not necessarily a uniform whole-montage effect.
Instead, the study is interested in whether abstract and concrete questions show a different spatial
pattern across the anterior versus posterior left lateral coverage.

\subsection{Core Analytical Decision}
The primary analysis in Step~4 should therefore be based on:
\begin{enumerate}
    \item \textbf{HbO only} as the primary chromophore;
    \item \textbf{long-separation channels only} for the inferential statistics;
    \item \textbf{two pre-specified ROIs} derived from the montage:
    \begin{itemize}
        \item left frontal-lateral ROI;
        \item left posterior-lateral ROI.
    \end{itemize}
    \item \textbf{one participant-level mean per condition and ROI};
    \item \textbf{one primary dissociation test} based on the difference of differences:
    \[
    \Delta_i =
    \bigl(\text{Front Abstract} - \text{Front Concrete}\bigr)_i
    -
    \bigl(\text{Back Abstract} - \text{Back Concrete}\bigr)_i.
    \]
\end{enumerate}

This primary dissociation score is mathematically equivalent to the
Condition $\times$ ROI interaction in a \(2 \times 2\) repeated-measures design,
but it is simpler to implement, explain, and report.

\subsection{Important Interpretation Rule}
This step can test for statistical significance, but it cannot guarantee significance.
The correct purpose of Step~4 is:
\begin{quote}
to test the abstract-versus-concrete difference cleanly and report honestly whether the data show
a statistically significant effect in the pre-specified front-versus-back left lateral comparison.
\end{quote}

\subsection{Inputs}
Step~4 begins only from the outputs of Step~3:
\begin{itemize}
    \item preprocessed participant-session fNIRS files;
    \item harmonized trigger table;
    \item session-level quality-control table;
    \item session problem log;
    \item preprocessing parameter file.
\end{itemize}

It also uses the outputs of Step~2:
\begin{itemize}
    \item question metadata;
    \item question type labels (\texttt{abstract} or \texttt{concrete});
    \item trial reconstruction from PsychoPy;
    \item question start and answer timing.
\end{itemize}

\subsection{Fixed Channel Layout}
All preprocessed participant-session FIF files share the same MNE channel layout and order.
The preprocessed data contain 60 channels in total:
\begin{itemize}
    \item 30 HbO channels;
    \item 30 HbR channels.
\end{itemize}

The exact channel labels are fixed across all preprocessed files.
This is important because it allows one consistent ROI definition for the whole dataset.

\subsection{Short-Channel Rule}
The study includes short-separation measurements, and these short channels must \textbf{not}
be used in the inferential ROI analysis.

The eight short source-detector pairs are:
\begin{verbatim}
S1_D8
S2_D9
S3_D10
S4_D11
S5_D12
S6_D13
S7_D14
S8_D15
\end{verbatim}

Therefore:
\begin{itemize}
    \item short channels are used only for nuisance correction and preprocessing in Step~3;
    \item short channels are excluded from all Step~4 ROI summaries;
    \item only long-separation channels are used for the inferential tests.
\end{itemize}

\subsection{Long-Channel Structure Used in Step~4}
After excluding the eight short pairs, the remaining long-separation pairs are:
\begin{verbatim}
S1_D3
S1_D6
S2_D4
S2_D6
S3_D1
S3_D3
S3_D7
S4_D2
S4_D4
S4_D7
S5_D3
S5_D4
S5_D6
S5_D7
S6_D1
S6_D2
S6_D5
S6_D7
S7_D2
S7_D5
S8_D1
S8_D5
\end{verbatim}

These 22 long-separation source-detector pairs form the full cortical long-channel set available
for analysis.

\subsection{ROI Definition Strategy}
The ROIs are defined from the montage at the channel level, using source-detector pairs.
To keep the primary test clean, the ROIs are divided into:
\begin{itemize}
    \item \textbf{front ROI channels};
    \item \textbf{back ROI channels};
    \item \textbf{transition channels}, which are excluded from the primary confirmatory analysis.
\end{itemize}

\subsection{Primary ROI Specification}

\subsubsection{Left Frontal-Lateral ROI}
The left frontal-lateral ROI is defined by the following HbO channels:

\begin{table}[ht]
\centering
\caption{Exact primary HbO channels for the left frontal-lateral ROI.}
\label{tab:front_roi_hbo}
\begin{tabular}{lll}
\hline
MNE channel name & Source-detector pair & Approximate topographic span \\
\hline
\texttt{S2\_D6 hbo} & S2--D6 & F7--F5 \\
\texttt{S2\_D4 hbo} & S2--D4 & F7--FT7 \\
\texttt{S1\_D6 hbo} & S1--D6 & F3--F5 \\
\texttt{S1\_D3 hbo} & S1--D3 & F3--FC3 \\
\texttt{S5\_D6 hbo} & S5--D6 & FC5--F5 \\
\texttt{S5\_D4 hbo} & S5--D4 & FC5--FT7 \\
\texttt{S5\_D3 hbo} & S5--D3 & FC5--FC3 \\
\hline
\end{tabular}
\end{table}

\subsubsection{Left Posterior-Lateral ROI}
The left posterior-lateral ROI is defined by the following HbO channels:

\begin{table}[ht]
\centering
\caption{Exact primary HbO channels for the left posterior-lateral ROI.}
\label{tab:back_roi_hbo}
\begin{tabular}{lll}
\hline
MNE channel name & Source-detector pair & Approximate topographic span \\
\hline
\texttt{S4\_D2 hbo} & S4--D2 & T7--TP7 \\
\texttt{S3\_D1 hbo} & S3--D1 & C3--CP3 \\
\texttt{S6\_D2 hbo} & S6--D2 & CP5--TP7 \\
\texttt{S6\_D1 hbo} & S6--D1 & CP5--CP3 \\
\texttt{S6\_D5 hbo} & S6--D5 & CP5--P5 \\
\texttt{S7\_D2 hbo} & S7--D2 & P7--TP7 \\
\texttt{S7\_D5 hbo} & S7--D5 & P7--P5 \\
\texttt{S8\_D1 hbo} & S8--D1 & P3--CP3 \\
\texttt{S8\_D5 hbo} & S8--D5 & P3--P5 \\
\hline
\end{tabular}
\end{table}

\subsubsection{Transition Channels Excluded from the Primary Confirmatory Test}
The following long HbO channels are treated as transition channels and excluded from the primary
confirmatory front-versus-back analysis:

\begin{table}[ht]
\centering
\caption{Transition HbO channels excluded from the primary confirmatory test.}
\label{tab:transition_hbo}
\begin{tabular}{lll}
\hline
MNE channel name & Source-detector pair & Approximate topographic span \\
\hline
\texttt{S5\_D7 hbo} & S5--D7 & FC5--C5 \\
\texttt{S4\_D4 hbo} & S4--D4 & T7--FT7 \\
\texttt{S4\_D7 hbo} & S4--D7 & T7--C5 \\
\texttt{S3\_D3 hbo} & S3--D3 & C3--FC3 \\
\texttt{S3\_D7 hbo} & S3--D7 & C3--C5 \\
\texttt{S6\_D7 hbo} & S6--D7 & CP5--C5 \\
\hline
\end{tabular}
\end{table}

\subsection{HbR Channel Mapping}
For the secondary HbR analysis, the same source-detector memberships should be used, replacing
the HbO suffix by the HbR suffix.
For example:
\begin{itemize}
    \item \texttt{S2\_D6 hbo} in the front ROI corresponds to \texttt{S2\_D6 hbr} in the HbR analysis;
    \item \texttt{S4\_D2 hbo} in the back ROI corresponds to \texttt{S4\_D2 hbr} in the HbR analysis.
\end{itemize}

Thus, the HbR ROI definitions are identical at the pair level.

\subsection{Primary Outcome}
The primary outcome should be the baseline-corrected mean HbO signal during the actual question window:
\begin{verbatim}
question_start --> answer
\end{verbatim}

This should be computed separately for:
\begin{itemize}
    \item each trial;
    \item each long channel;
    \item each ROI;
    \item each condition.
\end{itemize}

\subsection{Why This Level of Simplicity Is Appropriate}
The duration of a trial is not fixed.
Different participants answered different questions and took different amounts of time.
For a basic and transparent first inferential analysis, the cleanest strategy is:
\begin{enumerate}
    \item use the actual trial window from question start to answer;
    \item compute a simple trial-level fNIRS summary over that interval;
    \item average trial summaries within participant and condition;
    \item compare conditions with a within-subject test.
\end{enumerate}

This avoids overcomplicating the workflow while respecting the natural timing of the task.

\subsection{Participant Inclusion Rules}
A participant-session should be included in the \textbf{primary analysis} only if:
\begin{itemize}
    \item the final Step~3 status is \texttt{READY\_FOR\_STEP4} or
    \texttt{READY\_FOR\_STEP4\_WITH\_WARNINGS};
    \item the harmonized trigger table contains valid \texttt{question\_start} and \texttt{answer}
    events for the trials used;
    \item at least two valid abstract trials are available;
    \item at least two valid concrete trials are available;
    \item at least one valid front-ROI long HbO channel remains for that participant-session;
    \item at least one valid back-ROI long HbO channel remains for that participant-session.
\end{itemize}

Participants that fail these rules must remain listed in the exclusion table and must not be dropped silently.

\subsection{Trial Inclusion Rules}
A trial should be included if:
\begin{itemize}
    \item it has a valid \texttt{question\_start} trigger;
    \item it has a valid \texttt{answer} trigger;
    \item the question type is known (\texttt{abstract} or \texttt{concrete});
    \item the question window duration is positive;
    \item the preprocessed signal is available over that interval for at least one valid channel in each ROI.
\end{itemize}

A trial should be excluded if:
\begin{itemize}
    \item the start or answer trigger is unresolved;
    \item the condition label is missing;
    \item the signal is missing over the relevant interval;
    \item the baseline segment is unavailable;
    \item the trial belongs to a session marked \texttt{EXCLUDE\_FROM\_STEP4}.
\end{itemize}

\subsection{Definition of the Trial Window}
For each included trial:
\begin{itemize}
    \item the analysis window starts at the harmonized \texttt{question\_start};
    \item the analysis window ends at the harmonized \texttt{answer}.
\end{itemize}

This trial-specific interval is the primary task window for Step~4.

\subsection{Baseline Definition}
For each trial, define a simple baseline window immediately before question onset:
\begin{verbatim}
baseline window = [question_start - 2 s, question_start]
\end{verbatim}

The rule should be:
\begin{itemize}
    \item if the full baseline exists, use it;
    \item if the baseline is incomplete or missing, exclude that trial from the primary analysis;
    \item record the exclusion in the Step~4 trial log.
\end{itemize}

\subsection{Trial-Level Signal Summary}
For participant \(i\), channel \(j\), and trial \(k\), define the baseline-corrected trial response as:
\begin{equation}
r_{ijk}
=
\overline{x_{ijk}(t \in [t_{\mathrm{start}}, t_{\mathrm{answer}}])}
-
\overline{x_{ijk}(t \in [t_{\mathrm{start}}-2, t_{\mathrm{start}}])},
\end{equation}
where:
\begin{itemize}
    \item \(x_{ijk}\) is the preprocessed HbO signal;
    \item \(t_{\mathrm{start}}\) is the harmonized question start;
    \item \(t_{\mathrm{answer}}\) is the harmonized answer time.
\end{itemize}

\subsection{ROI-Level Trial Summaries}
For each trial, compute two ROI values.

\paragraph{Front ROI trial value}
\begin{equation}
y^{(F)}_{ik}
=
\frac{1}{J^{(F)}_{ik}}
\sum_{j \in F_{ik}} r_{ijk},
\end{equation}
where \(F_{ik}\) is the set of valid front-ROI HbO channels available for participant \(i\) on trial \(k\),
and \(J^{(F)}_{ik}\) is the number of those channels.

\paragraph{Back ROI trial value}
\begin{equation}
y^{(B)}_{ik}
=
\frac{1}{J^{(B)}_{ik}}
\sum_{j \in B_{ik}} r_{ijk},
\end{equation}
where \(B_{ik}\) is the set of valid back-ROI HbO channels available for participant \(i\) on trial \(k\),
and \(J^{(B)}_{ik}\) is the number of those channels.

If a channel is part of the ROI definition but marked bad or unavailable in a given participant-session,
it should simply be omitted from that participant-session's ROI average.
The number of valid channels contributing to each ROI must be recorded.

\subsection{Participant-Level Condition Means}
For each participant, compute the mean ROI response separately for abstract and concrete trials.

\paragraph{Front ROI}
\begin{equation}
\bar{y}^{(F,A)}_i
=
\frac{1}{n_i^{(A)}}
\sum_{k \in A} y^{(F)}_{ik},
\qquad
\bar{y}^{(F,C)}_i
=
\frac{1}{n_i^{(C)}}
\sum_{k \in C} y^{(F)}_{ik},
\end{equation}

\paragraph{Back ROI}
\begin{equation}
\bar{y}^{(B,A)}_i
=
\frac{1}{n_i^{(A)}}
\sum_{k \in A} y^{(B)}_{ik},
\qquad
\bar{y}^{(B,C)}_i
=
\frac{1}{n_i^{(C)}}
\sum_{k \in C} y^{(B)}_{ik},
\end{equation}
where:
\begin{itemize}
    \item \(A\) denotes abstract trials;
    \item \(C\) denotes concrete trials;
    \item \(n_i^{(A)}\) and \(n_i^{(C)}\) are the numbers of valid abstract and concrete trials.
\end{itemize}

\subsection{Primary Dissociation Score}
For each participant, define the primary dissociation score as:
\begin{equation}
\Delta_i
=
\bigl(\bar{y}^{(F,A)}_i - \bar{y}^{(F,C)}_i\bigr)
-
\bigl(\bar{y}^{(B,A)}_i - \bar{y}^{(B,C)}_i\bigr).
\end{equation}

This score answers the main Step~4 question directly:
does the abstract-versus-concrete contrast differ between the front and the back ROI?

\subsection{Primary Statistical Test}
The primary inferential analysis should be a \textbf{two-sided one-sample t-test} across participants on
the dissociation score \(\Delta_i\).

\paragraph{Null hypothesis}
\begin{equation}
H_0: \mu_{\Delta} = 0
\end{equation}

\paragraph{Alternative hypothesis}
\begin{equation}
H_1: \mu_{\Delta} \neq 0
\end{equation}

The significance threshold for the primary test should be:
\begin{verbatim}
alpha = 0.05
\end{verbatim}

\subsection{Why This Is the Right Primary Test}
This test is the cleanest primary choice because:
\begin{itemize}
    \item it directly targets the front-versus-back dissociation of interest;
    \item it stays fully within-subject;
    \item it is easy to interpret;
    \item it avoids unnecessary model complexity;
    \item it yields one main p-value for the primary hypothesis.
\end{itemize}

\subsection{Primary Reporting Quantities}
The main Step~4 results table must report:
\begin{itemize}
    \item number of included participants;
    \item mean front-abstract HbO response;
    \item mean front-concrete HbO response;
    \item mean back-abstract HbO response;
    \item mean back-concrete HbO response;
    \item mean dissociation score \(\Delta\);
    \item one-sample t statistic for \(\Delta\);
    \item degrees of freedom;
    \item exact p-value;
    \item 95\% confidence interval for the mean dissociation score;
    \item Cohen's \(d\) for the one-sample dissociation test.
\end{itemize}

\subsection{Planned Follow-Up Tests}
To describe the direction of the effect, Step~4 should also include two planned follow-up paired t-tests:

\paragraph{Front ROI follow-up}
\begin{equation}
H_0: \mu\bigl(\bar{y}^{(F,A)} - \bar{y}^{(F,C)}\bigr) = 0
\end{equation}

\paragraph{Back ROI follow-up}
\begin{equation}
H_0: \mu\bigl(\bar{y}^{(B,A)} - \bar{y}^{(B,C)}\bigr) = 0
\end{equation}

These follow-ups should:
\begin{itemize}
    \item be clearly labeled as planned follow-up tests;
    \item report exact p-values;
    \item report 95\% confidence intervals;
    \item report paired-effect sizes;
    \item be corrected across the two tests using a simple method such as Holm correction.
\end{itemize}

\subsection{Primary Conclusion Rule}
The final conclusion of Step~4 must be based first on the primary dissociation test.

The interpretation rule should be:
\begin{itemize}
    \item if the primary dissociation test is significant, conclude that the abstract-versus-concrete
    contrast differs significantly between the front and back left-lateral ROIs;
    \item if the primary dissociation test is not significant, conclude that this primary analysis does
    not show statistically significant evidence for a front-versus-back dissociation between abstract
    and concrete questions.
\end{itemize}

\subsection{Interpretation of the Sign of the Dissociation Score}
The sign of \(\Delta\) should be interpreted as follows:
\begin{itemize}
    \item if \(\Delta > 0\), the abstract-versus-concrete difference is more positive in the front ROI
    than in the back ROI;
    \item if \(\Delta < 0\), the abstract-versus-concrete difference is more positive in the back ROI
    than in the front ROI.
\end{itemize}

The planned follow-up tests are then used to describe whether this pattern reflects:
\begin{itemize}
    \item stronger abstract responses in the front ROI;
    \item stronger concrete responses in the back ROI;
    \item or both.
\end{itemize}

\subsection{Conclusion Templates}
The report should automatically generate one short conclusion paragraph.

\paragraph{Template if the primary dissociation test is significant.}
The pre-specified primary analysis showed a statistically significant front-versus-back dissociation
between abstract and concrete questions
(\(t(df)=\ldots\), \(p=\ldots\), mean dissociation \(=\ldots\), 95\% CI \([\ldots,\ldots]\), \(d=\ldots\)).
This indicates that the abstract-versus-concrete contrast was significantly different in the front and
back left-lateral ROIs.
Follow-up comparisons showed that \ldots

\paragraph{Template if the primary dissociation test is not significant.}
The pre-specified primary analysis did not show a statistically significant front-versus-back dissociation
between abstract and concrete questions
(\(t(df)=\ldots\), \(p=\ldots\), mean dissociation \(=\ldots\), 95\% CI \([\ldots,\ldots]\), \(d=\ldots\)).
In this dataset, the abstract-versus-concrete contrast was not significantly different between the
front and back left-lateral ROIs.

\subsection{Secondary Analyses}
To keep Step~4 simple, only a small number of secondary analyses should be included.

\subsubsection{Secondary Analysis 1: HbR ROI analysis}
Repeat the same ROI-based participant-level analysis using HbR instead of HbO.
This result must be clearly labeled as \textbf{secondary}.

\subsubsection{Secondary Analysis 2: Exploratory channel-level tests}
For each valid long channel in the left montage, compute the participant-level abstract-versus-concrete
contrast and run a paired comparison.
Because this introduces multiple comparisons, the resulting p-values must be corrected.
These analyses should be labeled \textbf{exploratory}.

\subsubsection{Secondary Analysis 3: Descriptive time-course plots}
Generate descriptive mean waveforms aligned to question onset for:
\begin{itemize}
    \item abstract trials;
    \item concrete trials;
    \item front ROI;
    \item back ROI.
\end{itemize}

These plots are useful for interpretation and quality control, but they do not replace the primary test.

\subsection{Analyses Intentionally Deferred to Step~5}
The following should be postponed to Step~5:
\begin{itemize}
    \item custom-duration GLM analysis;
    \item mixed-effects models with participant and item random effects;
    \item interactions with correctness;
    \item interactions with response time;
    \item interactions with question length;
    \item more detailed topographic or connectivity analyses.
\end{itemize}

\subsection{Expected Outputs}

\subsubsection*{Tables}
\begin{verbatim}
01_step4_inclusion_table.csv
02_trial_fnirs_summary_hbo.csv
03_trial_fnirs_summary_hbr.csv
04_participant_roi_condition_summary.csv
05_primary_dissociation_test_results.csv
06_followup_roi_tests.csv
07_secondary_hbr_results.csv
08_channel_exploratory_results.csv
09_step4_exclusion_log.csv
10_step4_final_conclusion.csv
\end{verbatim}

\subsubsection*{Figures}
\begin{verbatim}
figures/step4/front_back_fourcell_plot.png
figures/step4/primary_dissociation_distribution.png
figures/step4/front_roi_abstract_vs_concrete.png
figures/step4/back_roi_abstract_vs_concrete.png
figures/step4/front_hbo_waveform.png
figures/step4/back_hbo_waveform.png
figures/step4/channel_exploratory_map.png
\end{verbatim}

\subsubsection*{Reports}
\begin{verbatim}
reports/step4/step4_primary_analysis_report.tex
reports/step4/tables/primary_dissociation_table.tex
reports/step4/tables/followup_roi_tests_table.tex
reports/step4/text/final_conclusion.tex
\end{verbatim}

\subsection{Minimal Figures for the Main Report}
The main Step~4 report should include at least:
\begin{enumerate}
    \item a four-cell figure showing participant-level means for:
    Front--Abstract, Front--Concrete, Back--Abstract, Back--Concrete;
    \item a distribution plot of the participant-level dissociation scores \(\Delta_i\);
    \item one paired comparison plot for the front ROI;
    \item one paired comparison plot for the back ROI;
    \item descriptive HbO waveforms by condition for the front and back ROIs.
\end{enumerate}

\subsection{Quality-Control Checks Before Finalizing Step~4}
Before Step~4 is considered complete, verify that:
\begin{enumerate}
    \item all included sessions passed Step~3;
    \item every included trial has a valid condition label;
    \item every included trial has valid start and answer triggers;
    \item every included trial has a valid baseline segment;
    \item only long-separation channels are used in the inferential statistics;
    \item the eight short channel pairs are excluded from both primary ROIs;
    \item the front ROI is defined identically for all participants;
    \item the back ROI is defined identically for all participants;
    \item the six transition long channels are excluded from the primary confirmatory test;
    \item the primary dissociation test is run once and clearly labeled as primary;
    \item follow-up tests are clearly labeled as planned follow-ups;
    \item exact p-values appear in the final report.
\end{enumerate}

\subsection{Minimal Completion Criteria for Step~4}
Step~4 is complete when:
\begin{itemize}
    \item the ROI definitions have been frozen;
    \item the trial-level HbO summary table exists;
    \item the participant-level ROI-by-condition summary table exists;
    \item the primary dissociation test has been run;
    \item the primary p-value is reported;
    \item the planned follow-up contrasts have been reported;
    \item a short final conclusion has been generated;
    \item all exclusions and warnings are documented.
\end{itemize}

\subsection{Short Practical Summary}
In simple terms, Step~4 should:
\begin{quote}
use only long-channel HbO data, define a front left ROI and a back left ROI from the exact MNE channel
layout, summarize each trial from question start to answer, average those trial summaries within participant
for abstract and concrete questions, compute one front-versus-back dissociation score per participant,
test that score against zero, report the p-value, and write the conclusion clearly.
\end{quote}

\end{document}
